\section*{Bab \ref{ch:g8:combustion-and-fuels} --- Pembakaran: Bahan Bakar, Produk, dan Gas Rumah Kaca}

\begin{solution}{exo:g8:combustion-and-fuels:1}
\ce{C + O2 -> CO2}.
\end{solution}

\begin{solution}{exo:g8:combustion-and-fuels:2}
Dengan oksigen yang cukup, semua karbon menjadi karbon dioksida
(sempurna). Dengan oksigen yang terlalu sedikit, sebagian menjadi karbon
monoksida atau jelaga (tidak sempurna).
\end{solution}

\begin{solution}{exo:g8:combustion-and-fuels:3}
Gas itu tidak berwarna, tidak berbau, dan tidak berasa: tidak ada yang
menyadarinya, padahal gas itu beracun.
\end{solution}

\begin{solution}{exo:g8:combustion-and-fuels:4}
Fosil: batu bara, gas alam, bensin. Terbarukan: kayu, biogas, etanol
dari bit gula.
\end{solution}

\begin{solution}{exo:g8:combustion-and-fuels:5}
\ce{C3H8 + 5O2 -> 3CO2 + 4H2O}.
\end{solution}

\begin{solution}{exo:g8:combustion-and-fuels:6}
Karbon: 2 \ce{CO2}; hidrogen: 3 \ce{H2O}; oksigen di kanan $4 + 3 =
7$, yang 1 di antaranya berasal dari etanol, jadi 3 \ce{O2}:
\ce{C2H6O + 3O2 -> 2CO2 + 3H2O}.
\end{solution}

\begin{solution}{exo:g8:combustion-and-fuels:7}
Jelaga, yaitu karbon. \omterm{def:g4:what-a-fire-needs:burning}{Pembakarannya} tidak sempurna: nyalanya kekurangan
oksigen.
\end{solution}

\begin{solution}{exo:g8:combustion-and-fuels:8}
Sekitar 339 ppm pada tahun 1980 dan sekitar 370 ppm pada tahun 2000:
naik sekitar 30 ppm.
\end{solution}

\begin{solution}{exo:g8:combustion-and-fuels:9}
Karbon dalam kayu diambil dari \omterm{def:g4:air-a-mixture-of-gases:air}{udara} oleh pohon beberapa tahun atau
dasawarsa sebelumnya: membakarnya mengembalikan karbon dioksida itu.
Karbon dalam batu bara tersimpan di bawah tanah selama jutaan tahun:
membakarnya menambah karbon dioksida baru ke \omterm{def:g4:air-a-mixture-of-gases:air}{udara}.
\end{solution}

\begin{solution}{exo:g8:combustion-and-fuels:10}
\ce{2C + O2 -> 2CO}.
\end{solution}

\begin{solution}{exo:g8:combustion-and-fuels:11}
$(427 - 316)/316 = 111/316 \approx 0.35$: sekitar \qty{35}{\percent}.
\end{solution}

\begin{solution}{exo:g8:combustion-and-fuels:12}
Mula-mula oksigennya cukup: \ce{CH4 + 2O2 -> CO2 + 2H2O}. Ketika
oksigen di ruangan tertutup itu terpakai, nyalanya kekurangan \omterm{def:g4:air-a-mixture-of-gases:air}{udara}:
\ce{2CH4 + 3O2 -> 2CO + 4H2O}. Karbon monoksida, yang tidak berbau dan
beracun, menumpuk di dalam ruangan.
\end{solution}

\begin{solution}{pb:g8:combustion-and-fuels:1}
\textbf{1.} \ce{C3H8 + 5O2 -> 3CO2 + 4H2O}.

\textbf{2.} Kiri: 3 C, 8 H, 10 O. Kanan: 3 C, $4 \times 2 = 8$ H,
$3 \times 2 + 4 = 10$ O. Setara.

\textbf{3.} \ce{2C3H8 + 7O2 -> 6CO + 8H2O}.

\textbf{4.} \omterm{def:g8:combustion-and-fuels:complete}{Pembakaran sempurna}: 5 \omterm{def:g7:atoms-and-molecules:molecule}{molekul} oksigen per \omterm{def:g7:atoms-and-molecules:molecule}{molekul} propana,
dibandingkan $7/2 = 3.5$ untuk \omterm{def:g8:combustion-and-fuels:complete}{pembakaran tidak sempurna}.

\textbf{5.} \omterm{def:g4:air-a-mixture-of-gases:air}{Udara} di dalam tenda tertutup segera kekurangan oksigen;
nyalanya tidak lagi mendapat cukup oksigen dan \omterm{def:g4:what-a-fire-needs:burning}{terbakar} tidak sempurna,
berwarna kuning.

\textbf{6.} Karbon monoksida: tidak berwarna, tidak berbau (dan tidak
berasa).

\textbf{7.} GHS02, nyala api: sangat mudah menyala. GHS04, tabung gas:
gas yang disimpan di bawah tekanan.

\textbf{8.} Misalnya: pakailah hanya di \omterm{def:g4:air-a-mixture-of-gases:air}{udara} terbuka, jangan pernah di
dalam tenda atau ruangan tertutup; jauhkan dari apa pun yang dapat
\omterm{def:g4:what-a-fire-needs:burning}{terbakar}; ganti tabungnya jauh dari nyala api apa pun.

\textbf{9.} $132 \div 44 = 3$ kali lebih berat.

\textbf{10.} Dari oksigen \omterm{def:g4:air-a-mixture-of-gases:air}{udara}, yang berakhir di dalam karbon dioksida
(dan air).

\textbf{11.} \omterm{def:g8:combustion-and-fuels:fossil}{Bahan bakar fosil} (berasal dari gas alam atau minyak
bumi): membakarnya menambah karbon dioksida di \omterm{def:g4:air-a-mixture-of-gases:air}{udara}.

\textbf{12.} $450 \times 3 = \qty{1350}{g} = \qty{1.35}{kg}$ karbon
dioksida.
\end{solution}
